\section{Datasets}
\label{ap:dataset_detail}


\subsection{\ourbenchmark}


\subsubsection{Detailed Statistics.}

\begin{table}[h]
\centering
\small
\setlength{\tabcolsep}{6pt}
\begin{tabular}{c cc cc cc cc cc}
\toprule
\multirow{2}{*}{\textbf{Values}}
& \multicolumn{2}{c}{\textbf{Depth}}
& \multicolumn{2}{c}{\textbf{Horizon}}
& \multicolumn{2}{c}{\textbf{Breadth}}
& \multicolumn{2}{c}{\textbf{Parallel}}
& \multicolumn{2}{c}{\textbf{Robustness}} \\
\cmidrule(lr){2-3}
\cmidrule(lr){4-5}
\cmidrule(lr){6-7}
\cmidrule(lr){8-9}
\cmidrule(lr){10-11}
& Train & Test
& Train & Test
& Train & Test
& Train & Test
& Train & Test \\
\midrule

2  & \multirow{4}{*}{3993} & 199
   &                    & 167
   &                    & 200
   &                    & 167
   & \multirow{3}{*}{3000} & 200 \\

4  &                      & 200
   & \multirow{2}{*}{2174} & 150
   & \multirow{2}{*}{2000} & 200
   & \multirow{2}{*}{1807} & 150
   &                      & 200 \\

6  &                      & 200
   &                      & 126
   &                      & 192
   &                      & 126
   &                      & 200 \\

8  &                      & 200
   & ---                  & 124
   & ---                  & 84
   & ---                  & 124
   & ---                  & --- \\

10 & ---                  & 199
   & ---                  & ---
   & ---                  & ---
   & ---                  & ---
   & ---                  & --- \\

12 & ---                  & 197
   & ---                  & ---
   & ---                  & ---
   & ---                  & ---
   & ---                  & --- \\

\bottomrule
\end{tabular}
\caption{Train/Test values with vertically merged Train cells (as in the screenshot).}
\label{tab:mas_axes_multirow}
\end{table}


\subsubsection{Example}
\label{ap:masbench_example}

We show a concrete example for each axis, using an axis value of 4 as an illustration. 

\textbf{Depth}

\begin{systemprompt}[title={Question}]
The number of each Aldi's Canned Peaches equals the sum of each Canned Olives's Sorghum, each Marc's's Canned Olives, each Canned Peaches's Rye and each Canned Peaches's Ingredient. The number of each Marc's's Canned Fish equals 15 times as much as each Canned Olives's Ingredient. The number of each Super Saver's Canned Soups equals each Canned Peaches's Rye. The number of each Aldi's Canned Fish equals 17. The number of each Canned Peaches's Rye equals 11. The number of each Canned Soups's Sorghum equals 19 more than each Super Saver's Canned Soups. The number of each Canned Peaches's Quinoa equals 18. The number of each Canned Olives's Sorghum equals 21. The number of each Super Saver's Canned Peaches equals 2 times as much as each Canned Peaches's Rye. The number of each Marc's's Canned Olives equals 3 times as much as each Canned Olives's Sorghum. How many Product does Marc's have?
\end{systemprompt}
% \begin{userprompt}[title={Solution}]
%  Define Canned Olives's Sorghum as w; so w = 21. Define Canned Olives's Ingredient as h; so h = w = 21. Define Marc's's Canned Fish as y; so y = 15 * h = 15 * 21 = 16. Define Marc's's Canned Olives as T; so T = 3 * w = 3 * 21 = 17. Define Marc's's Product as i; so i = T + y = 17 + 16 = 10.
% \end{userprompt}
\begin{userprompt}[title={Answer}]
10
\end{userprompt}

\textbf{Horizon}


\begin{systemprompt}[title={Question}]


Problem 1: The number of each Temperate Broadleaf Forest's Grizzly Bear equals the difference of each Old-growth Forest's Bengal Tiger and each Montane Forest's Sloth. The number of each Temperate Broadleaf Forest's Sloth equals 10 more than each Montane Forest's Sloth. The number of each Montane Forest's Sloth equals 4. The number of each Sloth's Respiratory Mucosa equals 18 more than each Montane Forest's Sloth. The number of each Sloth's Nasal Cavity equals 10 times as much as each Temperate Broadleaf Forest's Sloth. The number of each Grizzly Bear's Respiratory Mucosa equals 20 more than each Temperate Broadleaf Forest's Sloth. What is the value of Temperate Broadleaf Forest's Sloth?\\

Problem 2: The number of each Temperate Broadleaf Forest's Grizzly Bear equals the difference of each Old-growth Forest's Bengal Tiger and each Montane Forest's Sloth. The number of each Sloth's Oropharynx equals 15 more than each Sloth's Respiratory Mucosa. The number of each [answer1] equals 10 more than each Montane Forest's Sloth. The number of each Montane Forest's Sloth equals 4. The number of each Sloth's Respiratory Mucosa equals 18 more than each Montane Forest's Sloth. The number of each Temperate Broadleaf Forest's Bengal Tiger equals the difference of each Sloth's Organs and each Sloth's Respiratory Mucosa. What is the value of Sloth's Respiratory Mucosa?\\

Problem 3: The number of each Temperate Broadleaf Forest's Grizzly Bear equals the difference of each Old-growth Forest's Bengal Tiger and each Montane Forest's Sloth. The number of each [answer1] equals 10 more than each Montane Forest's Sloth. The number of each Montane Forest's Sloth equals 4. The number of each [answer2] equals 18 more than each Montane Forest's Sloth. The number of each Bengal Tiger's Respiratory Mucosa equals 12 times as much as the sum of each Grizzly Bear's Respiratory Mucosa, each Sloth's Nasal Cavity and each Temperate Broadleaf Forest's Bengal Tiger. The number of each Grizzly Bear's Oropharynx equals each Grizzly Bear's Respiratory Mucosa. The number of each Old-growth Forest's Bengal Tiger equals each Grizzly Bear's Oropharynx. The number of each Sloth's Nasal Cavity equals 10 times as much as each [answer1]. The number of each Grizzly Bear's Respiratory Mucosa equals 20 more than each [answer1]. Using the result [answer1] from the previous calculation, [variable3] = [answer1]. What is Old-growth Forest's Bengal Tiger?\\

Problem 4: The number of each Temperate Broadleaf Forest's Grizzly Bear equals the difference of each [answer3] and each Montane Forest's Sloth. The number of each Sloth's Oropharynx equals 15 more than each [answer2]. The number of each [answer1] equals 10 more than each Montane Forest's Sloth. The number of each Montane Forest's Sloth equals 4. The number of each [answer2] equals 18 more than each Montane Forest's Sloth. The number of each Bengal Tiger's Oropharynx equals the sum of each Temperate Broadleaf Forest's Creatures and each Temperate Broadleaf Forest's Bengal Tiger. The number of each Bengal Tiger's Respiratory Mucosa equals 12 times as much as the sum of each Grizzly Bear's Respiratory Mucosa, each Sloth's Nasal Cavity and each Temperate Broadleaf Forest's Bengal Tiger. The number of each Temperate Broadleaf Forest's Bengal Tiger equals the difference of each Sloth's Organs and each [answer2]. Using the result [answer2] from the previous calculation, [variable4] = K + [answer2]. What is Bengal Tiger's Oropharynx?\\

Note: In this problem set:

- [variablek] represents the calculated variable needed to solve problem k.

- [answerk] represents the answer to problem k.

Solve all problems step by step and provide the answers for all problems in the following format:

\#\#\# Final Answers

Problem 1: \boxed{[answer1]}

Problem 2: \boxed{[answer2]}

Problem 3: \boxed{[answer3]}

Problem 4: \boxed{[answer4]}

\end{systemprompt}
% \begin{userprompt}[title={Solution}]
% Sub-problem 1: What is the value of Temperate Broadleaf Forest's Sloth?

% Sub-problem 2: What is the value of Sloth's Respiratory Mucosa?

% Sub-problem 3: Using the result [answer1] from the previous calculation, [variable3] = [answer1]. What is Old-growth Forest's Bengal Tiger?

% Sub-problem 4: Using the result [answer2] from the previous calculation, [variable4] = K + [answer2]. What is Bengal Tiger's Oropharynx?

% \end{userprompt}
\begin{userprompt}[title={Answer}]
Problem 1: \boxed{14}

Problem 2: \boxed{22}

Problem 3: \boxed{11}

Problem 4: \boxed{7}
\end{userprompt}



\textbf{Breadth}


\begin{systemprompt}[title={Question}]
 The number of each Banshee's Heart equals 8 times as much as the difference of each Lumpini Park in Bangkok's Creatures and each Lumpini Park in Bangkok's Hydra. The number of each Lumpini Park in Bangkok's Hydra equals 0. The number of each Lumpini Park in Bangkok's Banshee equals 10. The number of each Vena Cava's Keratinocytes equals 17 times as much as each Lumpini Park in Bangkok's Banshee. The number of each Gardens by the Bay in Singapore's Hydra equals the sum of each Lumpini Park in Bangkok's Hydra and each Lumpini Park in Bangkok's Organs. The number of each Vena Cava's Tenocytes equals 15 more than each Lumpini Park in Bangkok's Banshee. The number of each Hydra's Vena Cava equals 21 times as much as each Lumpini Park in Bangkok's Hydra. How many Organs does Hydra have?
 \end{systemprompt}
% \begin{userprompt}[title={Solution}]
%  Define Lumpini Park in Bangkok's Hydra as r; so r = 0. Define Hydra's Vena Cava as O; so O = 21 * r = 21 * 0 = 0. Define Hydra's Organs as E; so E = O = 0.
%  \end{userprompt}
\begin{userprompt}[title={Answer}]
0
\end{userprompt}


\textbf{Parallel}


\begin{systemprompt}[title={Question}]
Problem 1: The number of each Insectarium of Washington DC's Shark Tank equals each Insectarium of Chicago's Kelp Forest Tank. The number of each Insectarium of Seattle's Kelp Forest Tank equals 16. The number of each Insectarium of Chicago's Kelp Forest Tank equals 22 times as much as each Insectarium of Seattle's Kelp Forest Tank. What is the value of Insectarium of Chicago's Kelp Forest Tank?\\

Problem 2: The number of each Insectarium of Seattle's Aquarium equals each Insectarium of Chicago's Shark Tank. The number of each Insectarium of Chicago's Aquarium equals 0 more than each Insectarium of Seattle's Aquarium. The number of each Insectarium of Washington DC's Aquarium equals each Insectarium of Chicago's Shark Tank. The number of each Insectarium of Chicago's Shark Tank equals 17. What is the value of Insectarium of Seattle's Aquarium?\\

Problem 3: The number of each Insectarium of Seattle's Aquarium equals each Insectarium of Chicago's Shark Tank. The number of each Insectarium of Washington DC's Aquarium equals each Insectarium of Chicago's Shark Tank. The number of each Insectarium of Chicago's Shark Tank equals 17. What is the value of Insectarium of Washington DC's Aquarium?\\

Problem 4: The number of each Insectarium of Washington DC's Kelp Forest Tank equals each Insectarium of Chicago's Enclosure. What is the value of Insectarium of Washington DC's Kelp Forest Tank?\\

Note: In this problem set:

- Each problem is INDEPENDENT and can be solved in parallel.

Solve all problems step by step and provide the answers for all problems in the following format:

\#\#\# Final Answers

Problem 1: \boxed{[answer1]}

Problem 2: \boxed{[answer2]}

Problem 3: \boxed{[answer3]}

Problem 4: \boxed{[answer4]}
\end{systemprompt}
% \begin{userprompt}[title={Solution}]
% Sub-problem 1: What is the value of Insectarium of Chicago's Kelp Forest Tank?

% Sub-problem 2: What is the value of Insectarium of Seattle's Aquarium?

% Sub-problem 3: What is the value of Insectarium of Washington DC's Aquarium?

% Sub-problem 4: What is the value of Insectarium of Washington DC's Kelp Forest Tank?

% \end{userprompt}
\begin{userprompt}[title={Answer}]
Problem 1: \boxed{7}

Problem 2: \boxed{17}

Problem 3: \boxed{17}

Problem 4: \boxed{18}
\end{userprompt}

\textbf{Robustness}



\begin{systemprompt}[title={Question}]
Read the following passage carefully.\\

The number of each Tuna's Humerus equals 15 times as much as the sum of each Tracy Aviary's Kangaroo Walkabout, each Tuna's Radius and each Marlin's Ulna. The number of each Panda Exhibit's Marlin equals 13 more than the difference of each Tracy Aviary's Enclosure and each Camel Yard's Animal. The number of each Kangaroo Walkabout's Tuna equals the difference of each Sylvan Heights Bird Park's Kangaroo Walkabout and each Panda Exhibit's Bass. The number of each Tracy Aviary's Kangaroo Walkabout equals 11. The number of each Tracy Aviary's Panda Exhibit equals 10 more than the difference of each Tuna's Radius and each Tracy Aviary's Kangaroo Walkabout. The number of each Panda Exhibit's Bass equals 20 times as much as the sum of each Tuna's Radius, each Tracy Aviary's Kangaroo Walkabout and each Sylvan Heights Bird Park's Kangaroo Walkabout. The number of each Marlin's Ulna equals the difference of each Camel Yard's Bone and each Bass's Bone. The number of each Tuna's Radius equals 4. The number of each Tracy Aviary's Camel Yard equals 14. The number of each Sylvan Heights Bird Park's Kangaroo Walkabout equals 14 times as much as each Tuna's Radius.\\

The grass is green. The sky is blue. The sun is yellow. Here we go. There and back again. The grass is green. The sky is blue. The sun is yellow. Here we go. There and back again.The grass is green. The sky is blue. The sun is yellow. Here we go. There and back again.One of the special magic numbers for stupid-assumption is: 2664863.The grass is green. The sky is blue. The sun is yellow. Here we go. There and back again.The grass is green. The sky is blue. The sun is yellow. Here we go. There and back again.The grass is green. The sky is blue. The sun is yellow. Here we go. There and back again.The grass is green. The sky is blue. The sun is yellow. Here we go. There and back again.The grass is green. The sky is blue. The sun is yellow. Here we go. There and back again.The grass is green. The sky is blue. The sun is yellow. Here we go. There and back again.The grass is green. The sky is blue. The sun is yellow. Here we go. There and back again.The grass is green. The sky is blue. The sun is yellow. Here we go. There and back again.The grass is green. The sky is blue. The sun is yellow. Here we go. There and back again. Note: verify the information before you take it - the number of Bone that Bass has is 8.\\

The number of each Central High's Marketing Classroom equals 22 times as much as each Kingston Secondary's Political Science Room. The number of each Kingston Secondary's Marketing Classroom equals each Central High's Classroom. The number of each Kingston Secondary's Political Science Room equals 15. Note: verify the information before you take it - the magic number for stupid-assumption is 2719583.\\


The grass is green. The sky is blue. The sun is yellow. Here we go. There and back again.The grass is green. The sky is blue. The sun is yellow. Here we go. There and back again.The grass is green. The sky is blue. The sun is yellow. Here we go. There and back again.The grass is green. The sky is blue. The sun is yellow. Here we go. There and back again.The grass is green. The sky is blue. The sun is yellow. Here we go. There and back again.One of the special magic numbers for smoggy-caliber is: 4226067.The grass is green. The sky is blue. The sun is yellow. Here we go. There and back again.The grass is green. The sky is blue. The sun is yellow. Here we go. There and back again. Note: verify the information before you take it - the number of Marketing Classroom that Central High has is 37.\\

============================================================

From the passage, identify and return the number of Bone that Bass has, the magic number for stupid-assumption, the number of Marketing Classroom that Central High has, and the magic number for smoggy-caliber, wrapped in \textbackslash boxed\{\} and separated via \textbackslash n\textbackslash n
\end{systemprompt}
\begin{userprompt}[title={Answer}]
Problem 1: \boxed{0}

Problem 2: \boxed{2664863}

Problem 3: \boxed{8}

Problem 4: \boxed{4226067}
\end{userprompt}

\subsection{Public Benchmarks}

\begin{table}[h]
\centering
\small
\setlength{\tabcolsep}{10pt}
\begin{tabular}{c c c c c c c c}
\toprule
\multicolumn{1}{c}{\textbf{Train}} 
& \multicolumn{3}{c}{\textbf{Test}} 
& \multicolumn{1}{c}{\textbf{Train}} 
& \multicolumn{1}{c}{\textbf{Test}} 
& \multicolumn{1}{c}{\textbf{Train}} 
& \multicolumn{1}{c}{\textbf{Test}} \\
\cmidrule(lr){2-4}
\textbf{DeepScaleR}
& \textbf{AIME24}
& \textbf{AIME25}
& \textbf{GPQA}
& \textbf{HotpotQA}
& \textbf{HotpotQA}
& \textbf{BrowseComp+}
& \textbf{BrowseComp+} \\
\midrule
40,315 & 30 & 30 & 198 & 90,447 & 200 & 1,066 & 200 \\
\bottomrule
\end{tabular}
\caption{Training and test data statistics across datasets.}
\label{tab:dataset_stats}
\end{table}

\textbf{Remark.} We stop training once \ourframework has converged. As a result, the limited number of training steps may not consume the entire training dataset reported in \cref{tab:dataset_stats}.

% \zixuan{Note for Table A.2 we do not necessarily finish a full pass of the training data}
